\songtitle{\textchinese{兩種心動} (Liǎng zhǒng xīndòng)}

\tag*{2026}\tag{mandopop}\tag{mandarin-lyrics}
% -> generos/pop.tex?
\begin{notes}
Lyrics by Claude following specific instructions on genre and narrative elements: two messages at two in the morning, two kinds of heartbeat. Retrieved from Dvigitt's Suno page (V6-MINI, 14 November 2026).
\end{notes}
\begin{style}
Mandopop ballad, layered chorus harmonies widening over a warm, bittersweet mix; mid-tempo pulse; soft piano opening builds into a bright, full-band chorus, drops to piano and vocals for the stripped-back bridge, then rises in the key-changed final chorus; North Mandarin female vocals in a low, dark lower register, smooth and silky.
\end{style}
\begin{lyrics}
\songpart{Verse 1}
\begin{chinese}
凌晨兩點 手機亮了又暗\\
一則是你 說今晚月色很淡\\
一則是他 問我有沒有吃晚餐\\
我握著手機 像握著兩種答案
\end{chinese}

\songpart{Pre-Chorus}
\begin{chinese}
和你在一起 時間總是慢下來\\
我們共撐一把傘 話題永遠說不完\\
和他在一起 世界忽然變明亮\\
他笑著牽我的手 說別怕 往前
\end{chinese}

\songpart{Chorus}
\begin{chinese}
我愛你的 無話不談\\
也貪戀他 給的溫暖\\
心像兩條 平行的線\\
為什麼偏偏 在今夜相見\\
誰來告訴我 愛情的答案
\end{chinese}

\songpart{Verse 2}
\begin{chinese}
你懂我沉默時 沒說出口的話\\
他會逗我笑 讓我忘了牽掛\\
朋友說 愛情不必想太複雜\\
可是心裡 兩扇門 都在敲打
\end{chinese}

\songpart{Bridge (stripped back, piano and vocals)}
\begin{chinese}
也許不必 急著寫下答案\\
也許先學會 誠實面對自己\\
愛不是考題 沒有標準答案\\
我只想 好好愛 不辜負誰
\end{chinese}

\songpart{Final Chorus (key change up)}
\begin{chinese}
我還在學 怎麼選擇\\
但我知道 不想再逃避\\
無論最後 牽起誰的手\\
都要是 全心全意的我
\end{chinese}
\end{lyrics}

